\chapter{Introduction}

\section{Problem}

This dissertation studies a practical question at the intersection of
cryptography and verifiable computation: how can encrypted computation be
efficiently verified inside a zkVM?

The challenge is not only performance. It is also a trust question. In a normal
deployment, users often trust the algorithm specification and implementation
maintainers, even if they do not fully trust the machine running the workload.
In verifiable-computation settings, this decomposition becomes explicit: proofs
can reduce trust in the executor, but they do not remove the need to trust the
primitive choice, implementation quality, and software supply chain. The same
trust-in-tooling concern applies to the build, library, and proving-stack chain
that produces a zkVM guest binary.

In a zkVM, every guest instruction contributes to proving cost, so performance
assumptions from normal CPUs do not directly carry over.

\section{Why This Is Hard}

The problem is difficult for a second reason: naive performance transfer between
models often fails. Circuit-oriented literature focuses on multiplicative
complexity \cite{albrecht2015ciphers,jaques2019grover}, while RISC Zero's
execution model ties proving cost to cycles, memory behavior, and execution
traces. A primitive that is ``ZK-friendly'' in circuit form can still underperform
in VM form if its software path is instruction-heavy or matrix-intensive.

Naive implementation choices also fail for security reasons. Encrypting inside a
guest program does not automatically authenticate the returned transcript, and a
valid zkVM receipt does not by itself identify the sender of a proof artifact.
The system therefore has to separate three questions that are often conflated:
whether ciphertext hides plaintext, whether the ciphertext and tag were computed
by the fixed guest image, and whether an external party should be trusted as the
sender of that transcript.

Finally, empirical claims are hard to make repeatable. Proving jobs are long,
machine state affects timings, and artifact provenance must be preserved well
enough for later reruns. This is why the dissertation treats benchmark harness
design, run IDs, and result interpretation as part of the research method rather
than as incidental engineering.

\section{Motivation}

The motivation is to combine two requirements that are usually treated
separately: trustless execution, where a verifier checks correctness without
re-running the full computation, and confidentiality, where sensitive data
remains hidden while the computation is still auditable.

In this project, that combination is implemented in the RISC Zero host/guest
model. The host prepares inputs, invokes proving, and verifies receipts, while
the guest executes deterministic cryptographic logic and commits public outputs
to the journal. This split matters because moving logic between host and guest
changes both security scope and proving cost.

The resulting workflow is relevant for secure outsourced computation,
reproducible cryptographic experiments, and systems where verifiability is
required without exposing plaintext inputs.

One concrete example is privacy-preserving healthcare analytics: a provider can
use a zkVM guest to process sensitive cohort inputs and provide an auditor with
receipt-backed evidence that the fixed program executed correctly while only the
intended public outputs are disclosed.

The scope is intentionally bounded. This report studies one zkVM stack (RISC
Zero), uses benchmark artifacts for empirical claims, and presents a
structured game-based security sketch rather than a full mechanized proof.

\section{Aim}

The project has three concrete aims:

\begin{enumerate}
\item Benchmark cryptographic primitives inside the RISC Zero zkVM.
\item Identify which primitive and implementation strategy is most suitable for
this execution model.
\item Build an authenticated encryption pipeline that runs inside the zkVM and is
accompanied by a security argument.
\end{enumerate}

\section{Measurable Objectives}

The aims are evaluated through the following measurable objectives:

\begin{table}[t]
\centering
\scriptsize
\begin{tabular}{p{0.12\linewidth}p{0.39\linewidth}p{0.39\linewidth}}
\hline
ID & Objective & Evidence criterion \\
\hline
O1 & Implement comparable zkVM benchmark targets for AES, Salsa20, LowMC, and
operation-level kernels & Each enabled target emits structured JSON with
success status, timings, cycles, and benchmark ID provenance \\
O2 & Measure primitive and mode-family performance under one fixed RISC Zero
harness & Run IDs and aggregate files support all reported headline
tables and figures \\
O3 & Explain why LowMC theory did not translate into the best zkVM result in this
implementation & Analysis links LowMC/AES measurements to operation shape,
trace expansion, and implementation overhead \\
O4 & Build and evaluate an authenticated AES pipeline & CTR-only and CTR+HMAC
target families run end-to-end with receipt verification and matched block-size
coverage \\
O5 & State a scoped security argument and limitations & Threat model,
assumptions, game sketches, limitations, and future work are explicitly tied to
the implemented transcript format \\
\hline
\end{tabular}
\caption{Measurable objectives used to evaluate project success.}
\label{tab:measurable-objectives}
\end{table}

\section{Key Results Preview}

The evaluation produced three headline results:

\begin{itemize}
\item AES is the practical primitive choice in this RISC Zero implementation:
\texttt{lowmc-r0-optimised} is about \(34.85\times\) slower than \texttt{aes-r0}
on proving time in the main benchmark run.
\item LowMC optimisation was still useful locally, reducing proving time by about
\(25.59\%\) relative to the LowMC baseline snapshot, but the remaining absolute
gap to AES is too large for final selection.
\item In the authenticated AES pipeline, CTR+HMAC consistently increases
user-cycle work relative to CTR-only at all tested block sizes, while
single-trial wall-clock deltas are noisy enough that they are interpreted as
directional rather than statistically conclusive.
\end{itemize}

These results provide a scoped answer to the dissertation question: the strongest
claim is not that AES is always better than LowMC in all proof systems, but that
measured VM-level behavior in this stack supports AES-CTR plus HMAC as the final
implemented path.

\section{Limitations Preview}

The main limitations are stated here so the results are not overclaimed. All
headline measurements use RISC Zero only, most final values come from one trial
per target, and some provenance data such as per-target commit hashes is not
embedded directly in the benchmark JSON schema. Chapter 12 explains the impact
of each limitation and Chapter 14 prioritizes mitigation work.

\section{Contributions}

The main contributions are a reproducible benchmark setup for comparing
cryptographic workloads in RISC Zero, baseline implementations for AES, Salsa20,
and LowMC, and optimization workspaces focused on LowMC internals and AES S-box
backends.

The dissertation also provides a comparative analysis of AES and LowMC under
zkVM constraints, showing where circuit-level intuition and VM-level behavior
diverge in practice.

This positioning is deliberately different from prior work that primarily
optimizes or evaluates primitives in circuit and MPC settings. Here, the
contribution is receipt-backed VM-level benchmarking and system-level
interpretation.

Finally, it presents an authenticated pipeline for the local benchmark/proof
workflow using AES-CTR with HMAC-SHA256-192 in an Encrypt-then-MAC flow inside
the guest program, along with a structured IND-CPA game where the adversary also
receives the tag and RISC Zero receipt.

\section{Dissertation Structure}

The table of contents provides full navigation. Briefly, Chapter 2 covers
background and theory, Chapter 3 defines the measurement methodology, and
Chapter 4 presents benchmark framing and result slots. Chapter 5 analyzes LowMC
versus AES behavior, and Chapter 6 records final cipher selection.

Chapter 7 focuses on authentication-layer reasoning and design decisions, while
Chapter 8 details the authenticated AES pipeline design. Chapter 9 maps these
design decisions to repository implementation targets. Chapter 10 presents the
security game sketch, Chapter 11 presents evaluation results, Chapter 12
discusses limitations, Chapter 13 covers legal/social/ethical/professional
considerations, Chapter 14 outlines future work, and Chapter 15 concludes.
